\subsection{Lecture alignment checks}
Lecture 11 defines truncation error as the one step defect divided by $h$.
For $u'=u$, $u(0)=1$, each row below starts from the exact initial state.
The numerical update is unprojected RK4 and the reference is the exponential.
\begin{center}
\begin{tabular}{rrrr}
$h$ & Defect & Defect divided by $h$ & Defect order \\
\hline
@@TABLE_ROWS@@
\end{tabular}
\end{center}
The fifth order defect yields fourth order truncation error per unit time.
At very small steps floating point cancellation would limit this comparison.

Lectures 12 and 13 use the frozen Jacobian
$A=\left(\begin{smallmatrix}-2&0\\400&-2000\end{smallmatrix}\right)$.
Its eigenvalues are $-2$ and $-2000$, with magnitude ratio $1000$.
For $h=0.001$ the fast scaled eigenvalue is $-2$, inside the RK4 region.
For $h=0.0015$ it is $-3$, outside. This constant coefficient verification
fixture demonstrates the stability restriction. It does not classify the
full CubeSat system as stiff or nonstiff.

Following Lecture 3 and the grid norm discussion in Lecture 14, convergence
records additionally store the time weighted norm
$\sqrt{\int_0^T e(t)^2\,dt}$ using trapezoidal quadrature, and the largest
sampled error. Scalar error and attitude angle error in radians are separate.
No norm mixes quaternion components and dimensional angular rates.
Time weighting includes the shortened final interval and prevents sample
count from changing the meaning of a trajectory norm. Endpoint error remains
the primary convergence plot. Sampled norms are not continuous error bounds.
